\documentclass[a4paper, 10pt]{article}

% ---------- Pacotes essenciais ----------
\usepackage[T1]{fontenc}
\usepackage[utf8]{inputenc}
\usepackage[a4paper, top=14mm, bottom=14mm, left=10mm, right=10mm]{geometry}
\usepackage{xcolor}
\usepackage{enumitem}
\usepackage{ragged2e}
\usepackage{setspace}
\usepackage{array}
\usepackage{tabularx}
\usepackage{microtype}
\usepackage{parskip}

% ---------- Fontes ----------
\usepackage{mathpazo}      % Palatino — elegante e similar à Calluna
\newcommand{\Calluna}{\rmfamily}
\newcommand{\Courier}{\ttfamily}

% ---------- Cor e hiperlinks ----------
\definecolor{meuazul}{RGB}{0, 82, 155}
\definecolor{cinzafundo}{RGB}{242, 242, 242}
\definecolor{cinzaborda}{RGB}{190, 190, 190}

\usepackage[
colorlinks=true,
urlcolor=meuazul,
linkcolor=meuazul
]{hyperref}

% ---------- Caixa de fundo com mdframed ----------
\usepackage[framemethod=tikz]{mdframed}

\mdfdefinestyle{cvbox}{
	backgroundcolor=cinzafundo,
	linecolor=cinzaborda,
	linewidth=0.7pt,
	roundcorner=3pt,
	innertopmargin=6pt,
	innerbottommargin=6pt,
	innerleftmargin=8pt,
	innerrightmargin=8pt,
	leftmargin=0pt,
	rightmargin=0pt,
	skipabove=4pt,
	skipbelow=4pt,
	nobreak=true
}

% ---------- Ajustes gerais ----------
\setlength{\parindent}{0pt}
\setlength{\parskip}{0pt}

\newlength{\lblwidth}
\setlength{\lblwidth}{2.5cm}

% sectionBlock: full-width box wrapping title and content
\newcommand{\sectionBlock}[2]{%
	\par\addvspace{3pt}%
	\noindent%
	\begin{mdframed}[style=cvbox]%
		\noindent%
		\begin{minipage}[t]{\lblwidth}%
			\vspace{0pt}%
			\setstretch{1.0}%
			\raggedright%
			\small\scshape\Calluna%
			#1%
		\end{minipage}%
		\hspace{6pt}%
		\begin{minipage}[t]{\dimexpr\linewidth-\lblwidth-6pt\relax}%
			\vspace{0pt}%
			\setlength{\parindent}{0pt}%
			\setlength{\parskip}{2pt}%
			#2%
		\end{minipage}%
	\end{mdframed}%
	\par\addvspace{1pt}%
}

\renewcommand{\section}[1]{%
	{\small\bfseries\scshape #1}%
}

% ---------- Content helpers ----------
\newcommand{\cvBlock}{%
	\justifying%
	\setlength{\parindent}{0pt}%
	\setlength{\parskip}{2pt}%
	\noindent\ignorespaces%
}

\newcommand{\cvEntry}[4]{%
	\noindent%
	\begin{tabular*}{\linewidth}{@{}l@{\extracolsep{\fill}}r@{}}
		\textbf{#1} & {\small #2} \\
		\emph{#3} & {\small #4}
	\end{tabular*}%
	\par\vspace{1pt}%
}

\newenvironment{cvItems}{%
	\begin{itemize}[
		leftmargin=1.2em,
		labelsep=0.4em,
		itemsep=0pt,
		topsep=1pt,
		parsep=0pt,
		partopsep=0pt
		]
	}{%
	\end{itemize}
}

\newcommand{\cvLine}[1]{%
	\noindent #1\par\vspace{1pt}%
}

% ============================================================
\begin{document}
	\pagestyle{empty}
	
	% -------------------- CABEÇALHO --------------------
	\noindent
	\begin{minipage}[t]{0.68\textwidth}
		{\Calluna\fontsize{20pt}{26pt}\selectfont\textsc{Sergio Cordero Calvimontes}}
	\end{minipage}%
	\hfill
	\begin{minipage}[t]{0.28\textwidth}
		\raggedleft
		\begin{tabular}{@{}r@{}}
			{\Calluna\small +55 24 99995-9556} \\
			{\Calluna\small \href{http://lattes.cnpq.br/1576495718722378}{CV Lattes}}
		\end{tabular}
	\end{minipage}
	
	\par\vspace{1mm}
	\noindent\rule{\textwidth}{0.4pt}\par
	\vspace{2mm}
	
	% -------------------- RESUMO --------------------
	\sectionBlock{
		\section{Resumo}
	}{
		\cvBlock
		\cvLine{Engenheiro Mecatrônico com mestrado \textit{stricto sensu} em Engenharia Mecânica
			e Modelagem Computacional. Possui experiência em simulação numérica, computação científica
			e desenvolvimento de ferramentas computacionais. Atuação principal voltada à modelagem multiescala,
			materiais compósitos e processos termomecânicos. Atuação secundária voltada para o ensino de disciplinas
			eletroeletrônicas.}
	}
	
	% -------------------- FORMAÇÃO E PROJETOS --------------------
	\sectionBlock{
		\section{Formação\\[2pt]e Projetos}
	}{
		\cvBlock
		
		\cvEntry
		{EEIMVR-MCCT, UFF}
		{Rio de Janeiro, Brasil}
		{Mestrado em Modelagem Computacional em Ciência e Tecnologia}
		{2022--2024}
		\begin{cvItems}
			\item \href{https://drive.google.com/file/d/1fEubSSA7-rU-PD5rEgxHuHVYjIQxUjx9/view?usp=sharing}
			{Procedimentos de homogeneização computacional para análise multiescala.}
		\end{cvItems}
		
		\cvEntry
		{EEIMVR-PGMEC, UFF}
		{Rio de Janeiro, Brasil}
		{Mestrado em Engenharia Mecânica}
		{2016--2018}
		\begin{cvItems}
			\item \href{https://drive.google.com/file/d/1VzLID_a9mXDkNc7g3mvkYC9unA0O9GqH/view?usp=sharing}
			{Simulação numérica de processos eletrotérmico-mecânicos em soldagem.}
		\end{cvItems}
		
		\cvEntry
		{EESC / USP, UPB / USP}
		{São Paulo, Brasil / Cochabamba, Bolívia}
		{Bacharelado em Engenharia Mecatrônica}
		{2008--2013}
		\begin{cvItems}
			\item \href{https://drive.google.com/file/d/1wrtvWqw3lqqCWRn2qyCXGSgvPKS-htit/view?usp=sharing}
			{Medição indireta de forças de corte em processos CNC.}
		\end{cvItems}
		
		\cvEntry
		{Universidade Estácio de Sá}
		{Angra dos Reis, Brasil}
		{Bacharelado em Ciência da Computação, modalidade EAD}
		{2025--atual}
		
		\cvEntry
		{FORMA BRASIL EDUCACIONAL, CNE/CEB}
		{Paraíba, Brasil}
		{Técnico em Mecatrônica}
		{2024}
	}
	
	% -------------------- PERSPECTIVA DE CARREIRA --------------------
	\sectionBlock{
		\section{Perspectiva\\[2pt]de Carreira}
	}{
		\cvBlock
		\cvLine{\emph{Pesquisa técnico-científica}: Consolidação de manuscritos e materiais de apoio
			para submissão de artigos em computação, modelagem de materiais e simulação numérica.}
		\cvLine{\emph{Desenvolvimento acadêmico}: Organização de portfólio técnico-científico para
			candidatura ao Doutorado em Modelagem Computacional no LNCC-MCTI, Petrópolis, RJ.}
		\cvLine{\emph{Atuação técnica}: Acompanhamento de oportunidades na área pública e de engenharia,
			incluindo cadastro de reserva na Itaipu Binacional, usina nuclear Angra 3 e analista naval na EMGEPRON.}
	}
	
	% -------------------- PUBLICAÇÕES --------------------
	\sectionBlock{
		\section{Publicações\\[2pt]Destacadas}
	}{
		\cvBlock
		
		\cvLine{Calvimontes, S.\,C.
			\textbf{Projeto elétrico industrial básico: fundamentos, dimensionamento e documentação técnica.}
			1.\textsuperscript{a} ed. e-book.
			Angra dos Reis: Editora Argus Kepheus, 2026.
			[\href{https://drive.google.com/file/d/1j6KBMzZnQOnWV1OI968__dmG1RZcZ-Jl/view?usp=sharing}{PDF}]
			[\href{https://zenodo.org/records/20804262}{Zenodo}]
			[\href{https://archive.org/details/calvimontes-s.-c.-projeto-eletrico-fundamentos-dimensionamento-e-documentacao-tecnica-2026-1ed-182p}{Internet Archive}]
			[\href{https://www.academia.edu/169042224/Projeto_el\%C3\%A9trico_industrial_b\%C3\%A1sico}{Academia.edu}]}
		
		\cvLine{Calvimontes, S.\,C. \textit{et al.}
			\textbf{Computational homogenization of linear viscoelastic multi-phase composites:
				Effective relaxation estimation in perovskite-based solar panels.}
			\emph{Mechanics of Advanced Materials and Structures},
			Vol.\,33, No.\,1, 2456105, 2026.
			[\href{https://doi.org/10.1080/15376494.2025.2456105}{DOI}]
			[\href{https://drive.google.com/file/d/1LJM0_6I7H_WkGn2Orc6aN3Xmn1YhA52J/view?usp=sharing}{PDF}]}
	}
	
	% -------------------- HABILIDADES --------------------
	\sectionBlock{
		\section{Habilidades}
	}{
		\cvBlock
		\cvLine{\emph{Idiomas}: Espanhol nativo; português fluente, Celpe-Bras: 7+; inglês avançado, IELTS: 7+.}
		\cvLine{\emph{Linguagens de programação}: {\Courier MATLAB/GNU-Octave, Python, C++, Arduino/ESP32}.}
		\cvLine{\emph{Ferramentas}: \LaTeX{} / TeXstudio, Git/\href{https://github.com/Argus-Kepheus}{GitHub}, MS Office, Inkscape.}
	}
	
	% -------------------- FORMAÇÃO COMPLEMENTAR --------------------
	\sectionBlock{
		\section{Formação\\[2pt]Complementar}
	}{
		\cvBlock
		\cvLine{\emph{Programação}: C; códigos G/M para CNC; microcontroladores PIC, Arduino e ESP32; prototipagem e desenho técnico-industrial com manufatura aditiva em ABS e outros materiais poliméricos.}
		\cvLine{\emph{Automação}: Ladder Logic para CLPs; comandos eletroeletrônicos e circuitos de acionamento com relés, contatores, botoeiras, sensores, chaves seletoras, comutadores e intertravamentos.}
		\cvLine{\emph{CAD/CAM/CAE}: SolidWorks/ OnShape/FreeCAD, AutoCAD/ZWCAD BricsCAD, MasterCAM/Vectric Aspire/ VCarve, DIALux/Relux, DEFORM-3D, ANSYS, Linkage.}
		\cvLine{\emph{Engenharia}: instalações elétricas: NBR\,5410, NBR\,5413, NBR\,5419 e NBR\,14039;
			análise de falhas e resistência de materiais; metalografia e previsão de propriedades de aço com IA.}
	}
	
\end{document}
